% ═══════════════════════════════════════════════════════════════
% ARBEITSBLATT: parecer y conocer (Lehrbuch-Stil TEST)
% ═══════════════════════════════════════════════════════════════
% Testversion: LB-Layout adaptiert für AB
% Farbe nur konservativ: Titel + Aufgabennummern
% ═══════════════════════════════════════════════════════════════

\documentclass[10pt, a4paper]{article}

\usepackage[utf8]{inputenc}
\usepackage[T1]{fontenc}
\usepackage[ngerman]{babel}

% --- Schrift: Helvetica (serifenlos) ---
\usepackage{helvet}
\renewcommand{\familydefault}{\sfdefault}
\usepackage{microtype}

\usepackage[margin=2cm, top=2cm, bottom=1.8cm]{geometry}
\usepackage{enumitem}
\usepackage{tabularx}
\usepackage{booktabs}
\usepackage{array}

\usepackage{fancyhdr}
\usepackage{lastpage}
\usepackage{needspace}

\usepackage{xcolor}
\usepackage{tcolorbox}
\tcbuselibrary{skins}

% ═══════════════════════════════════════════════════════════════
% FARBEN (konservativ)
% ═══════════════════════════════════════════════════════════════

\definecolor{akzent}{HTML}{c41e3a}    % Nur für Titel + Aufgabennummern

% ═══════════════════════════════════════════════════════════════
% VARIABLEN
% ═══════════════════════════════════════════════════════════════

\newcommand{\thema}{parecer y conocer}
\newcommand{\klasse}{9}
\newcommand{\maxpunkte}{26}

% ═══════════════════════════════════════════════════════════════
% KOPF- UND FUSSZEILE
% ═══════════════════════════════════════════════════════════════

\pagestyle{fancy}
\fancyhf{}
\renewcommand{\headrulewidth}{0.4pt}
\renewcommand{\footrulewidth}{0pt}
\lhead{\footnotesize\textsc{Spanisch} | Klasse \klasse}
\rhead{\footnotesize Arbeitsblatt}
\cfoot{\footnotesize Seite \thepage/\pageref{LastPage}}

% ═══════════════════════════════════════════════════════════════
% BOXEN (Lehrbuch-Stil, S/W)
% ═══════════════════════════════════════════════════════════════

% Merkkasten (dicker Rahmen, weiß)
\newtcolorbox{merkkasten}[1][]{
    colback=white,
    colframe=black,
    boxrule=1.2pt,
    arc=0pt,
    left=8pt, right=8pt, top=6pt, bottom=6pt,
    fonttitle=\bfseries,
    title=#1
}

% Hinweis-Box (grauer Hintergrund)
\newtcolorbox{hinweisbox}{
    colback=black!5,
    colframe=black!40,
    boxrule=0.4pt,
    arc=0pt,
    left=6pt, right=6pt, top=4pt, bottom=4pt
}

% ═══════════════════════════════════════════════════════════════
% BEFEHLE
% ═══════════════════════════════════════════════════════════════

% Spanisch: fett (S/W)
\newcommand{\es}[1]{\textbf{#1}}

% Lücke
\newcommand{\luecke}[1]{\rule{#1}{0.4pt}}

% Punktebox (S/W)
\newcommand{\punktebox}[1]{\hfill\fbox{\small \_/#1}}

% Aufgabenumgebung (Nummer farbig als einziger Akzent)
\newenvironment{aufgabe}[1]{%
    \needspace{5\baselineskip}%
    \vspace{8pt}
    \noindent\textbf{\textcolor{akzent}{Aufgabe #1}}
    \vspace{4pt}
    \nopagebreak[4]%
}{%
    \vspace{4pt}
}

% Teilaufgaben
\newlist{teilaufgaben}{enumerate}{2}
\setlist[teilaufgaben,1]{label=\alph*), leftmargin=2em, itemsep=3pt}

\setlength{\parindent}{0pt}
\setlength{\parskip}{4pt}

% ═══════════════════════════════════════════════════════════════
% DOKUMENT
% ═══════════════════════════════════════════════════════════════

\begin{document}

% --- TITEL (einzige größere Farbnutzung) ---
\begin{center}
    {\large\bfseries\textcolor{akzent}{Arbeitsblatt: \thema}}
\end{center}

\vspace{2pt}

% --- Kopfzeile mit Name/Datum ---
\noindent
\begin{tabularx}{\textwidth}{@{}X X X@{}}
\textbf{Name:} \luecke{4cm} &
\textbf{Klasse:} \klasse &
\textbf{Datum:} \luecke{2.5cm}
\end{tabularx}

\vspace{1pt}
\hrule
\vspace{10pt}

% ═══════════════════════════════════════════════════════════════
% MERKKASTEN 1: parecer
% ═══════════════════════════════════════════════════════════════

\begin{merkkasten}[Merkkasten 1: parecer]
\es{parecer} = \luecke{3cm} / \luecke{3cm}

\vspace{4pt}

Konstruktion wie \es{gustar}:

\vspace{4pt}

\begin{tabular}{@{}ll@{}}
me parece / me parecen & = \luecke{4cm} \\[4pt]
te parece / te parecen & = \luecke{4cm} \\[4pt]
le parece / le parecen & = \luecke{4cm} \\
\end{tabular}

\vspace{4pt}

\textbf{Beispiel:} \es{¿Qué te parece Gonzalo?} -- \es{Me parece} \luecke{2cm}.
\end{merkkasten}

% ═══════════════════════════════════════════════════════════════
% AUFGABE 1
% ═══════════════════════════════════════════════════════════════

\begin{aufgabe}{1}
Ergänze mit der richtigen Form von \es{parecer}. \punktebox{6}
\end{aufgabe}

\begin{teilaufgaben}
    \item \es{¿Qué te} \luecke{2cm} \es{los deberes?} -- \es{Me} \luecke{2cm} \es{difíciles.}
    \item \es{¿Qué os} \luecke{2cm} \es{la película?} -- \es{Nos} \luecke{2cm} \es{muy divertida.}
    \item \es{A Clara le} \luecke{2cm} \es{que Gonzalo es un poco raro.}
    \item \es{¿Qué te} \luecke{2cm} \es{quedar a las ocho?} -- \es{Me} \luecke{2cm} \es{bien.}
\end{teilaufgaben}

% ═══════════════════════════════════════════════════════════════
% MERKKASTEN 2: conocer
% ═══════════════════════════════════════════════════════════════

\vspace{4pt}

\begin{merkkasten}[Merkkasten 2: conocer]
\es{conocer} = \luecke{3cm} / \luecke{3cm}

\vspace{4pt}

\textbf{Konjugation:}

\vspace{4pt}

\begin{tabular}{|c|c|c|c|c|c|}
\hline
yo & tú & él/ella & nosotros & vosotros & ellos \\
\hline
\luecke{1.3cm} & conoces & conoce & conocemos & conocéis & conocen \\
\hline
\end{tabular}

\vspace{6pt}

\textbf{WICHTIG:} \es{conocer} + \es{a} bei \luecke{2.5cm}!

\vspace{3pt}

\es{¿Conoces \textbf{a} María?} (Person) \hspace{8mm} \es{¿Conoces la ciudad?} (Sache)
\end{merkkasten}

% ═══════════════════════════════════════════════════════════════
% AUFGABE 2
% ═══════════════════════════════════════════════════════════════

\begin{aufgabe}{2}
Ergänze die richtige Form von \es{conocer}. \punktebox{7}
\end{aufgabe}

\begin{teilaufgaben}
    \item \es{Yo no} \luecke{2cm} \es{a tu hermano.}
    \item \es{¿Tú} \luecke{2cm} \es{la calle Trinidad?}
    \item \es{Nosotros} \luecke{2cm} \es{a toda tu familia.}
    \item \es{Mis hermanos} \luecke{2cm} \es{a Alfonso.}
    \item \es{Clara todavía no} \luecke{2cm} \es{muy bien a Nicolás.}
    \item \es{¿Vosotros} \luecke{2cm} \es{a los amigos de Gonzalo?}
    \item \es{Ado dice: ``Yo} \luecke{2cm} \es{todas las calles del centro.''}
\end{teilaufgaben}

% ═══════════════════════════════════════════════════════════════
% AUFGABE 3
% ═══════════════════════════════════════════════════════════════

\begin{aufgabe}{3}
\es{a} oder keine Präposition? Ergänze. \punktebox{6}
\end{aufgabe}

\begin{teilaufgaben}
    \item \es{¿Conoces} \luecke{1cm} \es{mi madre?}
    \item \es{No conozco} \luecke{1cm} \es{esta ciudad.}
    \item \es{Queremos conocer} \luecke{1cm} \es{tus amigos.}
    \item \es{¿Conocéis} \luecke{1cm} \es{el nuevo piso de Clara?}
    \item \es{Todavía no conozco} \luecke{1cm} \es{Nicolás.}
    \item \es{Mi hermano conoce} \luecke{1cm} \es{todas las canciones de este grupo.}
\end{teilaufgaben}

% ═══════════════════════════════════════════════════════════════
% AUFGABE 4
% ═══════════════════════════════════════════════════════════════

\begin{aufgabe}{4}
Schreibe eigene Sätze. \punktebox{7}
\end{aufgabe}

\textbf{a)} Beantworte die Fragen mit \es{parecer}: (3 P.)

\begin{itemize}[leftmargin=2em, itemsep=2pt]
    \item \es{¿Qué te parece tu clase de español?}

    \vspace{2pt}
    \luecke{12cm}

    \item \es{¿Qué te parecen los deberes de matemáticas?}

    \vspace{2pt}
    \luecke{12cm}
\end{itemize}

\textbf{b)} Schreibe 2 Sätze mit \es{conocer}. Achte auf \es{a} bei Personen! (4 P.)

\vspace{4pt}
1. (Person): \luecke{10cm}

\vspace{4pt}
2. (Sache): \luecke{10cm}

% ═══════════════════════════════════════════════════════════════
% PUNKTEÜBERSICHT
% ═══════════════════════════════════════════════════════════════

\vspace{12pt}
\hrule
\vspace{8pt}

\begin{flushright}
\textbf{Punkte:} \luecke{1cm} / \maxpunkte
\end{flushright}

\end{document}
